% SimGolf-Latent: NeurIPS style full paper draft
\documentclass{article}
\usepackage[preprint]{neurips_2024}
\usepackage[utf8]{inputenc}
\usepackage{amsmath,amsfonts,amssymb,graphicx,booktabs}
\usepackage[colorlinks=true,citecolor=blue,urlcolor=blue]{hyperref}
\title{SimGolf in Latent Space: Combining DreamerV3 World Models with Simulation-Based Global Search}
\author{Author Name \\
Affiliation \\
\texttt{email@institution.edu} \\
(Replace with your author list and affiliations before submission)}
\date{}

\begin{document}
\maketitle

\begin{abstract}
We present SimGolf-Latent, an approach that integrates simulation-based global
search (SimGolf) with learned latent world models (DreamerV3's RSSM). When the
agent detects high decision importance (large TD-error or model uncertainty),
it resets imagination to stored latent checkpoints and performs global branching
search in latent space. Branch returns are penalized by model uncertainty and
aggregated to yield TD(\(\lambda\)) targets and policy updates. Our
implementation supports CEM branch proposals, uncertainty penalties, and
branch-to-real distillation. We provide algorithmic details, code, and an
experimental protocol for hard navigation tasks.
\end{abstract}

\section{Introduction}
Model-based reinforcement learning (RL) has shown strong sample efficiency by
learning compact latent representations and performing optimization in latent
space. Dreamer-family agents use a Recurrent State-Space Model (RSSM) to
capture dynamics and perform actor-critic updates via imagined rollouts, which
dramatically reduces sample complexity in many domains. However, these
approaches still struggle on **long-horizon**, **sparse-reward** tasks that
require non-local planning and explicit backtracking: local greedy updates may
fail to discover long-range sequences of actions that lead to sparse rewards.

SimGolf \cite{sim_golf} proposed simulation-based global search on top of
episodic resets to rediscover and refine long-range plans; the central idea is
to perform branching search from previously visited states to explore
alternative continuations. In this work we adapt and extend SimGolf to operate
entirely in **latent space** using DreamerV3's RSSM. Latent resets eliminate
the need for costly environment resets while allowing broad, imagined
branching that can be used to refine actor and critic at critical decision
points.

Contributions: we make the following contributions in this paper:
\begin{enumerate}
\item We present SimGolf-Latent, a planner that performs global branching in the
latent space of a learned RSSM and integrates its imagined branches into
actor-critic training via TD(\(\lambda\)) targets and distillation.
\item We introduce a lightweight trigger module that decides when to invoke
expensive planning based on TD-errors, model uncertainty, and policy entropy,
with cooldowns to limit compute overhead.
\item We implement an efficient CEM-based branch proposal module that can be
warm-started from the policy and supports uncertainty-penalized returns.
\item We provide a robust open-source scaffold (planner utilities, training
script, examples, tests), reproducibility guidance, and a detailed experimental
protocol to evaluate the planner on hard navigation benchmarks.
\end{enumerate}

In the rest of this paper we formalize the planner, describe implementation
choices and integration with DreamerV3, and present a rigorous experimental
protocol and templates for reporting results.

\section{Method}
\subsection{Overview and notation}
Let \(z_t\in\mathbb{R}^d\) denote the learned stochastic latent state (RSSM). The
world model parameterized by \(\theta\) defines transitions
\(p_\theta(z_{t+1}\mid z_t,a_t)\), a decoder \(p_\theta(o_t\mid z_t)\), a reward head
\(r_\theta(z_t)\), and an optional uncertainty estimate \(\sigma_{r,\theta}(z_t)\)
for predicted rewards. The actor \(\pi_\phi(a\mid z)\) and critic
\(v_\psi(z)\) operate in latent space and are trained using imagined rollouts.

\subsection{Checkpoint buffer}
We maintain a fixed-capacity checkpoint buffer \(C\) storing tuples
\((z, s, t, m)\) where \(z\) is a latent, \(s\) is a numeric score (e.g., absolute TD-error),
\(t\) is the env step when recorded, and \(m\) are optional metadata (entropy,
novelty). The buffer supports:
\begin{itemize}
\item insertion with overwrite (FIFO when full),
\item prioritized sampling proportional to positive scores (falling back to uniform), and
\item save/load to disk for reproducibility.
\end{itemize}
Implementation: see `planner.checkpoint_buffer.CheckpointBuffer`.

\subsection{Trigger logic}
A lightweight trigger decides when global planning is worth the compute. Typical
signals include: large absolute TD-error, model ensemble variance on rewards,
and extreme policy entropy (either very low or very high). Triggers obey a
cooldown to prevent overuse. Formally, with thresholds \(\delta_{TD},\tau_{unc}\),
trigger when
\( |\delta_{TD}|>\delta_{TD} \) OR \( \mathrm{Var}(r)>\tau_{unc} \) OR
\( H(\pi(\cdot|z))\notin[H_{low},H_{high}]\). See `planner.triggers.TriggerConfig`.

\subsection{Branch generation and scoring}
From a saved latent \(z_s\), we generate \(B\) candidate branches each of
horizon \(H\) using either policy sampling or CEM-optimized action sequences.
During each imagined step the RSSM produces \(z_{t+1}, r_t, \gamma_t\) and
(optional) reward uncertainty \(\sigma_{r,t}\). We use an uncertainty-penalized
reward
\( r'_t = r_t - \kappa\sigma_{r,t} \), where \(\kappa\ge0\) is a hyperparameter.
Branch returns are bootstrapped with the critic as:
\begin{equation}
\hat R_b = \sum_{k=0}^{H-1}\Big(\prod_{j=0}^{k-1}\gamma_{j}\Big) r'_{k} + \Big(\prod_{j=0}^{H-1}\gamma_j\Big) v_\psi(z_{H}).
\end{equation}
Branches are ranked by \(\hat R_b\) and aggregated for learning or action
selection.

\subsection{Algorithm (pseudocode)}
Below is high-level pseudocode for the planner integration used in our
experiments. This pseudocode highlights the planner trigger, checkpoint sampling,
branch generation (policy sampling or CEM), branch scoring, and imagined
updates applied to the actor and critic.

\begin{verbatim}
# Main training loop (simplified)
for step in range(total_steps):
    # Real environment interaction and Dreamer updates (world model + actor/critic)
    o, a, r, done = env.step(a_real)
    z = rssm.observe(o, a, r)
    update_world_model_and_actor_critic(replay)

    # compute trigger metrics
    td_err = compute_td_error(z)
    unc = compute_model_uncertainty(z)
    ent = compute_policy_entropy(z)

    # maybe record checkpoint
    if should_record_checkpoint(td_err, unc, ent):
        buffer.push(z, score=td_err, step=step)

    # planner invocation
    if trigger(td_err, unc, ent) and len(buffer) > 0:
        z_saved = buffer.sample(k=1, prioritized=True)[0]['z']
        branches = simulate_branches(rssm, actor, value_fn, z_saved, cfg)
        # compute TD(lambda) targets and update actor/critic from imagined data
        imagined_updates = build_imagined_targets(branches, cfg.lambda)
        apply_imagined_updates(imagined_updates)
        # optional: distill top-k first actions into policy
        if cfg.branch_distill:
            actions, probs = branch_to_action_distribution(branches, cfg.topk, cfg.tau)
            distill_actions_into_policy(actions, probs)

    # choose next real action (may be biased by planner)
    a_real = select_real_action(z, actor, planner_hint)

# end loop
\end{verbatim}

We provide reference implementations in `planner/` and example notebooks to
visualize branches and their effects.

\subsection{CEM branch proposal}
When CEM is enabled we optimize an action-sequence distribution to maximize
predicted branch return. The implemented CEM uses a population of size
\(P\), selects the top-\(E\) elites each iteration, and re-fits a Gaussian
sequence model for several iterations. CEM is warm-started from policy samples
to improve sample efficiency and may be constrained to reasonable action bounds.
Implementation details and pragmatic defaults live in `planner.advanced`.

\subsection{Learning from branches}
Each branch yields an imagined trajectory \(\tau=(z_0,a_0,r_0,\gamma_0,\dots)\).
We compute TD(\(\lambda\)) targets backward along each branch using
Equation~\eqref{eq:td-lambda} and treat the results as additional imagined
training examples for the critic and actor. Concretely, for a single branch of
length \(H\) we compute targets via the backward recursion:
\begin{align}
G_{H} &= v_\psi(z_H),\nonumber \\
G_{t} &= r_t + \gamma_t\big((1-\lambda)v_\psi(z_{t+1}) + \lambda G_{t+1}\big)\quad\text{for }t=H-1,\dots,0.
\end{align}
The critic minimizes the MSE loss
\(L_{\rm critic}=\mathbb{E}_{\tau}[\|v_\psi(z_t)-\mathrm{stopgrad}(G_t)\|^2]\)
averaged over selected branches. The actor is updated using an advantage-weighted
policy gradient loss computed from imagined data:
\begin{equation}
L_\pi = -\mathbb{E}_{\tau, t}[\log\pi_\phi(a_t\mid z_t)\cdot\mathrm{stopgrad}(G_t - v_\psi(z_t))] - \eta\mathcal{H}(\pi_\phi).
\end{equation}
To compute stable gradient targets we stop gradients through \(G_t\) and use
bootstrapped critic values only where appropriate. Aggregation strategies
include using the top-\(k\) branches, normalizing branch weights with a
softmax of returns, or temperature-scaling to control variance introduced by
outlier branches.

\subsection{Branch-to-real distillation}
To bias real actions toward promising branches we extract the first-action
distribution over top branches and (optionally) distill that into the policy
via a supervised loss (KL or cross-entropy for discrete actions). `planner.distill`
provides helpers `branch_to_action_distribution` and `select_branch_action`.

\subsection{Computational complexity and practical tips}
Per planner invocation, the imagined-step cost scales as \(O(BH\cdot C_{step})\)
where \(C_{step}\) is the RSSM step cost. CEM adds a factor proportional to the
population size and iterations. To control wall-clock cost: reduce \(B\) or
\(H\), increase trigger thresholds, or limit CEM usage. Use mixed precision and
batch branch rollouts for GPU efficiency.

\section{Implementation details}
We provide a production-ready scaffold and practical guidance for integration.

\subsection{Code layout and key modules}
\begin{itemize}
\item `planner/checkpoint_buffer.py` — checkpoint storage, prioritized sampling, save/load.
\item `planner/simgolf_latent.py` — policy-sampled branching (minimal, robust interface).
\item `planner/advanced.py` — CEM branch proposal, uncertainty-penalized returns.
\item `planner/distill.py` — utilities for extracting and distilling first-action distributions.
\item `scripts/train_simgolf_dreamerv3.py` — integration scaffold showing how to:
  \begin{enumerate}
  \item compute trigger statistics (TD-error, uncertainty, entropy),
  \item sample checkpoints from buffer, and
  \item run `simulate_branches` and apply imagined updates.
  \end{enumerate}
\end{itemize}

\subsection{Integration hooks}
To integrate with an existing Dreamer agent:
\begin{enumerate}
\item instrument the training loop to compute TD-errors per latent and push to the checkpoint buffer when thresholds are met,
\item expose a periodic planner hook (or event-based trigger) that samples a checkpoint and runs the planner, and
\item provide a pathway to apply imagined branch targets to actor and critic updates (e.g., `DreamerAgent.update_actor_critic_from_branches`).
\end{enumerate}

\subsection{Logging schema and metrics}
Log planner-specific metrics alongside standard training metrics:
`planner/branches` (count), `planner/top_return`, `planner/trigger_rate`,
`planner/fps` (imagined steps per second), `planner/cost_fraction` (fraction of
wall-time used by planner). Also log `td_error`, `uncertainty`, and `entropy` to
facilitate ablation and analysis. Tests are available in `CA13/tests`.

\subsection{Developer notes}
Keep planner code import-safe (no heavy side-effects), add type hints and
clear docstrings, and add unit tests that cover edge cases (empty buffer,
file save/load, actor missing sample method). Use the existing stubs in the
training script for smoke-testing before wiring into full DreamerV3 agents.

\subsection{API signatures and expected shapes}
To ease integration we provide minimal recommended signatures and tensor shapes
for the key planner primitives (Python type hints shown informally):
\begin{itemize}
\item `simulate_branches(rssm, actor, value_fn, z_saved, cfg) -> List[Branch]`:
  - `z_saved`: Tensor `[1, z_dim]` or `[batch, z_dim]` (uses batch size 1 for checkpoints),
  - returns `Branch` objects with `.ret` (float) and `.traj` as list of `(z, a, r, gamma)`.
\item `CheckpointBuffer.push(z: Tensor, score: float, step: int)` and `sample(k:int, prioritized:bool)`:
  - `z` shape `[1, z_dim]`, `sample` returns list of dicts `{'z':Tensor,'score':float,'step':int}`.
\item `branch_to_action_distribution(branches, topk_frac, tau) -> (actions:Tensor[k,adim], probs:Tensor[k])`:
  - `actions` stacked first-actions for top-k branches, `probs` softmax-weighted on returns.
\item CEM helper: returns a list of `H` action tensors `[1, adim]` or a Tensor `[H,adim]` representing the proposed sequence.
\end{itemize}
These minimal contracts make it straightforward to plug in different RSSMs or
actor implementations without committing to a heavy dependency on a specific
project layout.

\section{Experimental setup and ablations}
We evaluate on hard navigation tasks where long-horizon planning and backtracking
matter. Below we provide experimental protocols to ensure fair comparison and
reproducibility.

\subsection{Benchmarks and environments}
Primary environments:
\begin{itemize}
\item D4RL AntMaze: `antmaze-umaze`, `antmaze-medium-play`, `antmaze-large-play` (use d4rl~v0.5+),
\item DMControl maze/nav variants (if available), and
\item Custom continuous mazes used for controlled ablations (deterministic map layouts).
\end{itemize}

\subsection{Baselines and variants}
\begin{itemize}
\item DreamerV3 baseline (no planner) using identical model capacity and optimization hyperparameters.
\item DreamerV3 + SimGolf-Latent (policy-sampled branches).
\item DreamerV3 + SimGolf-Latent + CEM (CEM enabled for branch proposals).
\end{itemize}

\subsection{Evaluation metrics and protocol}
\begin{itemize}
\item Success rate (binary goal completion) over evaluation episodes.
\item Episode return and steps-to-goal for successful episodes.
\item Planner overhead: imagined steps per second and wall-clock training time ratio.
\item Trigger statistics: average triggers per 1e5 steps and distribution across trajectories.
\end{itemize}

\subsection{Ablations}
Systematic ablations to understand behavior and sensitivity to planner
hyperparameters:
\begin{enumerate}
\item Branching factor \(B\): \{0 (off), 8, 16, 32\}.
\item Horizon \(H\): \{5, 10, 20\}.
\item Trigger modes: TD-error only, uncertainty only, entropy only, OR combination.
\item CEM vs policy sampling and CEM hyperparameters (population, elites, iters).
\item Uncertainty penalty \(\kappa\): \{0.0, 0.1, 0.5\}.
\item Buffer size: \{256, 1024, 4096\}.
\item Branch-to-real distillation: on/off and distillation temperature \(\tau\).
\end{enumerate}

\subsection{Evaluation protocol and defaults}
We recommend the following practical defaults and evaluation procedure to make
comparisons fair and reproducible (modify per environment):
\begin{table}[ht]
\centering
\begin{tabular}{lcc}
\toprule
Parameter & Default & Notes \\
\midrule
Branching factor (B) & 16 & planner branches per invocation \\
Horizon (H) & 12 & imagined steps per branch \\
Uncertainty penalty (\kappa) & 0.1 & reduces optimism from uncertain rewards \\
CEM pop & 64 & if CEM enabled \\
CEM elite & 8 & \\
CEM iters & 3 & \\
Buffer size & 1024 & checkpoint capacity \\
Trigger cooldown & 8 steps & prevents overuse \\
Eval frequency & 50k env steps & evaluate policy every N steps \\
Eval episodes & 100 & per seed per evaluation \\
Seeds & 5 & matched seeds across baselines \\
\bottomrule
\end{tabular}
\caption{Representative experimental defaults used in our protocol.}
\end{table}

For each experimental condition:
\begin{itemize}
\item run at least 5 seeds with matched RNGs, log per-step metrics, and save planner checkpoints,
\item perform evaluations every `Eval frequency` steps over `Eval episodes` episodes per seed, and
\item report mean ± std, IQM and 95% bootstrap confidence intervals for success rates and returns.
\end{itemize}

\subsection{Compute budgeting}
Report wall-clock training times and compute budgets (number of GPUs, GPU type,
CPU/RAM). For planner ablations also report planner fraction (fraction of
training wall-time consumed by planning). This helps explain trade-offs between
improved performance and added compute.
IQM where appropriate, and compute bootstrap confidence intervals for success
rates. Perform paired tests between baseline and planner variants when seeds
are matched.

\section{Results (Template)}
\label{sec:results}
The following templates and examples show how to present experimental results. Replace the example numbers with empirical values produced by your runs, and save all figures in the `figures/` directory. Use provided notebooks under `notebooks/` to generate plots and `scripts/` to aggregate metrics into CSVs for table construction.

\subsection{Tables}
\begin{itemize}
\item Final performance table: for each environment, report mean ± std success rate and mean episode return across seeds.
\item Overhead table: report planner imaginary steps per second and wall-clock training-time ratio vs baseline.
\end{itemize}

\begin{table}[ht]
\centering
\begin{tabular}{lccc}
\toprule
Method & AntMaze-Umid (success) & AntMaze-Med (success) & Wall-time overhead \\
\midrule
DreamerV3 (baseline) & 0.45 ± 0.12 & 0.38 ± 0.10 & 1.00x \\
DreamerV3 + SimGolf (B=16,H=12) & 0.62 ± 0.08 & 0.54 ± 0.09 & 1.22x \\
\bottomrule
\end{tabular}
\caption{Example: replace numbers with empirical results (mean ± std). Include number of seeds in caption.}
\end{table}

\subsection{Figures and curves}
\begin{itemize}
\item Success rate (y-axis) vs training steps (x-axis) with shaded episode-wise std across seeds.
\item Trigger rate over training (how often planner fired per 1e5 steps).
\item Ablation plots: performance metric vs \(B\) / \(H\) / \(\kappa\).
\end{itemize}

\subsection{Figures & artifacts}
Provide these concrete artifacts alongside the paper submission:
\begin{itemize}
\item `figures/success_curve.pdf/png` — success rate curves per env (SVG/PDF preferred for publication quality), 1200×600 px or vector.
\item `figures/trigger_rate.pdf` — trigger frequency vs steps.
\item `figures/ablation_B_H.pdf` — ablation grids for branching/horizon.
\item `figures/branch_value_distributions.pdf` — distributions of branch returns.
\item `data/metrics_<experiment>.csv` — step-indexed metrics for reproducibility and table generation.
\end{itemize}

Tips: prefer vector formats (PDF/SVG) for plots used in LaTeX. Include a small script `notebooks/generate_and_visualize_example.py --out_prefix figures/fig_branch_example` to recreate example figures from saved CSVs.

\subsection{Statistical reporting and analysis}
Report number of seeds, compute budget, and the random seeds used. For primary
comparisons, use paired tests when seeds are matched and report bootstrap CIs
for medians or means where appropriate. Discuss effect sizes and practical
significance, not only p-values.

\subsection{Hyperparameters (defaults)}
\begin{table}[ht]
\centering
\begin{tabular}{lcc}
\toprule
Parameter & Default & Notes \\
\midrule
Branching factor \(B\) & 16 & per planner call \\
Horizon \(H\) & 12 & imagined steps \\
CEM pop & 64 & if CEM enabled \\
CEM elite & 8 & \\
Kappa \(\kappa\) & 0.0 & uncertainty penalty \\
Trigger TD threshold & 0.7 & absolute TD-error \\
Buffer size & 1024 & \\
Actor lr & 3e-4 & Adam \\
Critic lr & 3e-4 & Adam \\
\bottomrule
\end{tabular}
\caption{Representative defaults (see configs/antmaze_simgolf.yaml).}
\end{table}

\section{Related work}
DreamerV1--V3 (Hafner et al.) introduced latent RSSMs and imagination-based policy
learning. SimGolf introduced episodic resets with branching search. CEM-MPC is a
classical planner; our CEM runs in latent-space to produce candidate action
sequences. Prioritized replay and contrastive representation learning inform our
checkpoint scoring and buffer design.

\section{Discussion, Limitations, and Future Work}
\subsection{Discussion}
The SimGolf-Latent planner is complementary to Dreamer-style imagination: it
injects global search at decision-critical moments, which can facilitate
backtracking and recovery in sparse-reward navigation tasks. Our design trades
extra compute (imagined rollouts) for better long-horizon planning and can be
tuned to the available budget via trigger thresholds, branching factor, and
horizon.

\subsection{Limitations}
Key limitations include:
\begin{itemize}
\item Model quality dependence: if the RSSM is biased or poorly calibrated, branch
    returns may be misleading and can degrade performance.
\item Computational overhead: frequent planning with large \(B,H\) or CEM can
    significantly increase wall-clock training time.
\item Partial observability and non-modeled hidden variables may make latent
    resets unreliable for certain tasks where important state is not captured.
\end{itemize}

\subsection{Future work}
Future directions to improve practicality and robustness:
\begin{itemize}
\item Better uncertainty estimation (ensembles or Bayesian heads) and principled
    ways to incorporate model epistemic uncertainty into branch scoring.
\item Cache-aware CEM warm-starts and hierarchical planner combinations to
    improve sample efficiency of CEM proposals.
\item Application to multi-agent or partially-observed multi-objective tasks
    where global resets could be coordinated or constrained.
\end{itemize}

\section{Conclusion}
We proposed SimGolf-Latent, a practical method for integrating global branching
search with Dreamer-style latent imagination. We provide code, tests, and a
clear experimental recipe to evaluate on hard navigation tasks and include an
extensible scaffold to reproduce experiments and iterate on planner design.

% -----------------------------
\appendix

\section{Worked Example: SimGolf-Latent Branching and TD($\lambda$) Calculation}
To illustrate the planner, we provide a toy example with $B=2$ branches, $H=3$ steps, and simple scalar rewards/discounts. Assume:
\begin{itemize}
\item $z_0$ is the checkpoint latent; $v_\psi(z)$ returns $0.5$ for all $z$.
\item Branch 1: $r = [1.0, 0.0, 0.0]$, $\gamma = [0.9, 0.9, 0.0]$
\item Branch 2: $r = [0.0, 0.0, 2.0]$, $\gamma = [0.9, 0.9, 0.0]$
\item $\lambda = 0.8$
\end{itemize}
For Branch 1, compute $G_3 = v_\psi(z_3) = 0.5$.
\begin{align*}
G_2 &= r_2 + \gamma_2 \big((1-\lambda)v_\psi(z_3) + \lambda G_3\big) = 0.0 + 0.0 \cdot ((1-0.8)\cdot 0.5 + 0.8 \cdot 0.5) = 0.0 \\
G_1 &= r_1 + \gamma_1 \big((1-\lambda)v_\psi(z_2) + \lambda G_2\big) = 0.0 + 0.9 \cdot (0.2 \cdot 0.5 + 0.8 \cdot 0.0) = 0.09 \\
G_0 &= r_0 + \gamma_0 \big((1-\lambda)v_\psi(z_1) + \lambda G_1\big) = 1.0 + 0.9 \cdot (0.2 \cdot 0.5 + 0.8 \cdot 0.09) \approx 1.0 + 0.9 \cdot (0.1 + 0.072) = 1.0 + 0.9 \cdot 0.172 = 1.0 + 0.1548 = 1.1548
\end{align*}
Branch 1's TD($\lambda$) targets: $[1.1548, 0.09, 0.0]$.

Branch 2 (similarly): $G_3 = 0.5$, $G_2 = 2.0$, $G_1 = 0.0 + 0.9 \cdot (0.2 \cdot 0.5 + 0.8 \cdot 2.0) = 0.9 \cdot (0.1 + 1.6) = 0.9 \cdot 1.7 = 1.53$, $G_0 = 0.0 + 0.9 \cdot (0.2 \cdot 0.5 + 0.8 \cdot 1.53) = 0.9 \cdot (0.1 + 1.224) = 0.9 \cdot 1.324 = 1.1916$.
Branch 2's TD($\lambda$) targets: $[1.1916, 1.53, 2.0]$.

This illustrates how the planner aggregates imagined returns and bootstraps with the critic.

\section{Glossary of Symbols}
\begin{tabular}{ll}
$z_t$ & RSSM latent at time $t$ \\
$a_t$ & action at time $t$ \\
$r_t$ & reward at time $t$ \\
$\gamma_t$ & discount at time $t$ \\
$v_\psi(z)$ & value function (critic) \\
$\pi_\phi(a|z)$ & policy (actor) \\
$B$ & branching factor (number of branches) \\
$H$ & branch horizon (steps per branch) \\
$\kappa$ & uncertainty penalty coefficient \\
$C$ & checkpoint buffer capacity \\
$\lambda$ & TD($\lambda$) parameter \\
\end{tabular}

\section{Best Practices for Using SimGolf-Latent}
\begin{itemize}
\item Tune $B$ and $H$ to match your compute budget; start with $B=8$, $H=10$.
\item Use prioritized checkpoint sampling for more effective planning.
\item Enable CEM only if policy sampling is insufficient; CEM is more expensive.
\item Log planner overhead and trigger frequency to avoid excessive compute use.
\item Always run ablations on $\kappa$ and buffer size to check robustness.
\item Use mixed precision and batch branch rollouts for GPU efficiency.
\end{itemize}

\section{Frequently Asked Questions (FAQ)}
\begin{itemize}
\item \textbf{Q: Can I use SimGolf-Latent with any world model?}
A: Yes, as long as your model exposes a step function and value head.
\item \textbf{Q: What if my actor does not have a sample method?}
A: The planner tries sample, act, forward, and __call__ in order.
\item \textbf{Q: How do I visualize branches?}
A: Use the provided notebook and scripts to plot branch returns and latent trajectories.
\item \textbf{Q: How do I know if the planner is helping?}
A: Compare success rates and wall-clock time to a DreamerV3 baseline; look for improved performance on sparse-reward tasks.
\item \textbf{Q: What if my buffer fills up?}
A: The buffer is FIFO; old checkpoints are overwritten. Increase capacity if you need more diversity.
\end{itemize}

% -----------------------------
\section{Appendix: README-to-Report Expanded Explanations}
This appendix provides a comprehensive, section-by-section academic expansion of the README.md, ensuring all concepts, algorithms, and practicalities are fully described for both research and tutorial use.

\subsection{Executive Summary}
SimGolf-Latent fuses simulation-based global search (SimGolf) with DreamerV3’s latent world model. The core idea is to maintain a buffer of high-importance latent checkpoints (e.g., high TD-error, high uncertainty, or near-goal), and, when triggered, reset the imagination to one of these checkpoints. The planner then performs branching rollouts in latent space, aggregates returns, and updates the policy/value functions. This approach enables broad exploration and replanning, especially in sparse-reward, long-horizon domains such as D4RL AntMaze and DMControl navigation.

\subsection{Selected Papers}
The method builds on two key lines of work: (1) SimGolf (NeurIPS 2024), which introduced simulation-based global search with episodic resets and branching, and (2) DreamerV3, which uses a learned RSSM for latent imagination and actor-critic updates. SimGolf-Latent combines these by performing SimGolf-style branching entirely in the latent space of DreamerV3.

\subsection{Novel Research Question}
The central question is: Can SimGolf’s global search be realized entirely in latent space, eliminating the need for real-environment resets while retaining exploration benefits? The hypothesis is that latent resets plus imagined branching reduce compounding error at critical decision points, improving success on hard navigation tasks.

\subsection{Core Idea and Pipeline}
The pipeline is as follows:
\begin{enumerate}
\item Collect real trajectories with DreamerV3’s actor.
\item Train RSSM, decoder, actor, and critic as in DreamerV3.
\item Compute uncertainty, TD-error, or novelty; push latent checkpoints into a buffer.
\item When triggered, pick a checkpoint, run SimGolf branching in latent for horizon $H$ and branching factor $B$.
\item Aggregate returns, update actor/critic from imagined branches, and optionally bias the next real action.
\item Repeat until the training budget is exhausted.
\end{enumerate}

\subsection{Mathematical Formulation}
\textbf{Latent Dynamics:} The RSSM transition is $p_\theta(z_{t+1} \mid z_t, a_t)$, with decoder $p_\theta(o_t \mid z_t)$, reward head $p_\theta(r_t \mid z_t)$, and discount head $p_\theta(\gamma_t \mid z_t)$.

\textbf{Actor-Critic:} The actor $\pi_\phi(a_t \mid z_t)$ is optimized on imagined rollouts, and the critic $v_\psi(z_t)$ is trained with TD($\lambda$) targets.

\textbf{SimGolf Branching:} For a checkpoint $z_s$, generate $B$ branches $\{\tau_b\}$, each a sequence $(z_s, a_{s,b,0:H-1}, z_{s,b,1:H})$. The branch return is:
\begin{equation}
R_b = \sum_{k=0}^{H-1} \left( \prod_{j=0}^{k-1} \gamma_{s,b,j} \right) r_{s,b,k} + \left( \prod_{j=0}^{H-1} \gamma_{s,b,j} \right) v_\psi(z_{s,b,H})
\end{equation}
Branches are aggregated (top-k or softmax) to update the policy/value.

\textbf{Trigger Conditions:} Triggers fire when $|v_\psi(z_t) - \text{target}_t| > \delta$, uncertainty is high, or the agent is near a goal.

\subsection{Algorithm Pseudocode}
See Section~\ref{sec:results} for a full pseudocode block. The key steps are: (1) collect real data, (2) update models, (3) compute triggers, (4) run planner if triggered, (5) update actor/critic from branches, (6) repeat.

\subsection{Architecture}
\textbf{Shared with DreamerV3:} RSSM, convolutional encoder/decoder, reward/discount/value heads.

\textbf{Additions for SimGolf-Latent:} (1) Checkpoint buffer for latents and metadata, (2) branching planner using RSSM and critic, (3) trigger module, (4) branch selection policy (top-1, mixture, or distillation).

\subsection{Losses}
\textbf{World Model Loss:} Reconstruction, reward, discount, and KL regularization.

\textbf{Actor Loss:} Maximize imagined return plus entropy regularization.

\textbf{Critic Loss:} TD($\lambda$) regression.

\textbf{Planner Auxiliary Loss:} Optional imitation/distillation of top-branch action distribution.

\subsection{Trigger Mechanisms}
Triggers can be based on TD-error, uncertainty, entropy, or periodic intervals. They are combined with OR logic and rate-limited by a cooldown.

\subsection{Hyperparameters}
Ablate over $B$ (branching factor), $H$ (horizon), buffer size, trigger thresholds, CEM settings, and distillation temperature. See Section~\ref{sec:results} for a full table.

\subsection{Benchmarks and Evaluation}
Evaluate on D4RL AntMaze, custom mazes, and DMControl navigation. Metrics include success rate, steps to goal, model loss, planner overhead, trigger frequency, and ablation results.

\subsection{Logging and Visualization}
Log all model and planner losses, planner stats (branches, horizon, FPS, trigger rate), and success/return metrics. Visualize success curves, trigger frequency, branch value distributions, and planning cost.

\subsection{Implementation Plan and Checklist}
Start from DreamerV3, add checkpoint buffer, trigger computation, latent branching, planner update path, logging, and ablations. See the checklist in the README for step-by-step progress.

\subsection{Reproducibility}
Seed all RNGs, log git hash and config YAML, save checkpoints and metrics, and provide exact configs for each experiment. Include commands for environment setup, training, testing, and paper build.

% -----------------------------
\section{Appendix: Code Walkthrough and API Usage}
This appendix provides a detailed walkthrough of the main code modules, classes, and functions in the SimGolf-Latent package, with explanations, integration tips, and example usage for each public API.

\subsection{planner/checkpoint	extunderscore buffer.py}
\textbf{Class: CheckpointBuffer}
\begin{itemize}
\item \texttt{__init__(capacity, device)}: Initializes a fixed-size buffer for storing latent checkpoints, scores, and step indices. The buffer supports prioritized sampling and is device-aware (CPU/GPU).
\item \texttt{push(z, score, step)}: Adds a new checkpoint (latent $z$) with a score (e.g., TD-error) and step index. Overwrites oldest if full.
\item \texttt{sample(k, prioritized)}: Samples $k$ checkpoints, either proportional to score (prioritized) or uniformly. Returns a list of dicts with $z$, score, step, and index.
\item \texttt{save(path)}, \texttt{load(path)}: Save/load buffer contents to disk (CPU tensors for portability).
\item \texttt{to(device)}: Moves all stored tensors to the specified device.
\end{itemize}
\textbf{Integration tip:} Use one buffer per agent or experiment; push latents when TD-error or uncertainty is high.

\subsection{planner/triggers.py}
\textbf{Class: TriggerConfig}
\begin{itemize}
\item Stores thresholds for TD-error, uncertainty, entropy, and cooldown period.
\end{itemize}
\textbf{Function: should	extunderscore trigger(td	extunderscore error, unc, entropy, cfg, last	extunderscore trigger, step)}
\begin{itemize}
\item Returns True if any trigger condition is met and cooldown has elapsed.
\end{itemize}
\textbf{Integration tip:} Call this function at each step to decide whether to invoke the planner.

\subsection{planner/simgolf	extunderscore latent.py}
\textbf{Function: simulate	extunderscore branches(rssm, actor, value	extunderscore fn, z	extunderscore saved, cfg)}
\begin{itemize}
\item Simulates $B$ branches of length $H$ from a checkpoint $z_{saved}$ using the provided RSSM, actor, and value function. Returns a list of Branch objects, each with a return and trajectory.
\item Handles actor methods (sample, act, forward, call) and RSSM step methods (step, imagine	extunderscore step, transition, predict).
\end{itemize}
\textbf{Integration tip:} Use this as the main planner entry point; pass in your DreamerV3 RSSM, actor, and critic.

\subsection{planner/advanced.py}
\textbf{Function: simulate	extunderscore branches	extunderscore advanced(...)}
\begin{itemize}
\item Adds CEM-based action sequence optimization and uncertainty-penalized returns. CEM is warm-started from the policy and can be tuned via population, elite count, and iterations.
\end{itemize}
\textbf{Integration tip:} Enable CEM in your config to use this advanced planner; otherwise, the default is policy sampling.

\subsection{planner/distill.py}
\textbf{Function: branch	extunderscore to	extunderscore action	extunderscore distribution(branches, topk	extunderscore frac, tau)}
\begin{itemize}
\item Extracts the first actions from the top-k branches and returns a softmax-weighted distribution over them (temperature $\tau$).
\end{itemize}
\textbf{Function: select	extunderscore branch	extunderscore action(branches, ...)}
\begin{itemize}
\item Samples a first action from the branch action distribution.
\end{itemize}
\textbf{Integration tip:} Use these to distill branch policies into your actor or to bias the next real action.

\subsection{scripts/train	extunderscore simgolf	extunderscore dreamerv3.py}
\textbf{Main script for running the planner in a DreamerV3-style training loop.}
\begin{itemize}
\item Loads config, sets up buffer, RSSM, actor, and value function (stubs or real models).
\item At each step: computes trigger metrics, pushes checkpoints, samples and runs planner if triggered, and prints planner stats.
\end{itemize}
\textbf{Example usage:}
\begin{verbatim}
python scripts/train_simgolf_dreamerv3.py --config configs/antmaze_simgolf.yaml --workdir ./runs/exp1
\end{verbatim}

\subsection{tests/}
\textbf{Unit tests for all major modules:}
\begin{itemize}
\item test	extunderscore checkpoint	extunderscore buffer.py: tests push, sample, save/load.
\item test	extunderscore simgolf	extunderscore latent.py: tests branch simulation and return ordering.
\item test	extunderscore advanced	extunderscore and	extunderscore distill.py: tests advanced planner and distillation utilities.
\item test	extunderscore triggers.py: tests trigger logic and cooldown.
\item test	extunderscore package	extunderscore api.py: tests package-level exports and API surface.
\end{itemize}
\textbf{Integration tip:} Run \texttt{pytest -q} in the CA13 directory to verify all functionality before running experiments.

\subsection{Example: End-to-End Planner Call}
\begin{verbatim}
from planner import CheckpointBuffer, simulate_branches, should_trigger, TriggerConfig

buffer = CheckpointBuffer(capacity=1024, device=device)
if should_trigger(td_err, unc, entropy, trig_cfg, last_trigger, step):
    sample = buffer.sample(k=1, prioritized=True)
    if sample:
        z_saved = sample[0]['z']
        branches = simulate_branches(rssm, actor, value_fn, z_saved, planner_cfg)
        a_plan = select_branch_action(branches)
        # use a_plan to bias next real action or update actor/critic
\end{verbatim}

% -----------------------------

\bibliographystyle{unsrt}
\begin{thebibliography}{9}
\bibitem{hafner2023dreamerv3}
Danijar Hafner et al., DreamerV3: Mastering Diverse Domains through World Models, ICLR 2023.
\bibitem{sim_golf}
SimGolf authors, Simulation-Based Global Search, NeurIPS 2024.
\end{thebibliography}

\end{document}
















